\subsection{四単純基点の元係数と全共有支持}
\label{subsec:cubic-CY-four-point-support}
\begin{lemma}\label{lem:cubic-CY-four-point-support}
零交換子CYの基点型1111では、印付き対称関係を
\[
 Q=z_0z_5+\lambda z_1z_4-(1+\lambda)z_2z_3,\qquad \lambda\ne0,-1
\]
とした全CY係数は、既存FS1--FS2の21混合係数と一つの交代二次係数 $g$ で尽くされる。
共有係数の非零支持は4096個、全monomial群の置換部分による軌道は218個であり、
次の互いに重ならない五場合へ漏れなく分かれる。
\[
\begin{array}{c|r|r}
\text{支持型}&\text{支持数}&S_4\text{軌道数}\\\hline
\text{全零}&1&1\\
\text{出次数1の頂点あり}&3471&173\\
\text{上記なし、出次数2の頂点あり}&609&40\\
\text{出次数0又は3、部分活性}&14&3\\
\text{十二辺全非零}&1&1
\end{array}
\]
出次数1の枝は $x_1\ne0,x_3=x_4=0$、
出次数2の枝は $x_1x_4\ne0,x_3=0$ へ移せ、各非零共有係数を1に規格化できる。
\end{lemma}
\begin{proof}
対称基底 $U$ と交代基底 $K$ を生成元対 $ab,ac,ad,bc,bd,cd$ の順に取る。
$A=\tfrac12\operatorname{Hess}(Q)$ とし、FS2の行列を明記すると
\[
 B=\begin{pmatrix}
 x_0&x_1&x_3&-x_6&-x_{10}&-x_{11}+x_8\\
 x_1&x_2&x_4&x_7&-x_{15}+x_8&-x_{16}\\
 x_3&x_4&x_5&x_{11}-x_{15}&x_{12}&x_{17}\\
 x_6&x_7&x_8&x_9&x_{13}&-x_{18}\\
 x_{10}&x_{11}&x_{12}&x_{13}&x_{14}&x_{19}\\
 x_{15}&x_{16}&x_{17}&x_{18}&x_{19}&x_{20}
 \end{pmatrix}.
\]
$C_{14}=C_{41}=(\lambda+1)/8$、$C_{23}=C_{32}=-\lambda/8$、他は零とする。
$J_{05}=J_{50}=1,J_{14}=J_{41}=-1,J_{23}=J_{32}=1$ とすると
\[
 UAU^{\mathsf T}+UBK^{\mathsf T}+KB^{\mathsf T}U^{\mathsf T}
       +K(C+gJ)K^{\mathsf T}
\]
は全係数で反巡回である。
混合36方向の巡回条件はrank15、交代対称21方向の条件はrank20、
合体した57方向の条件はrank35である。
従って表示した独立な21混合係数と $g$ が全解を与える。
必要なrank6条件は従来通り $C+gJ-B^{\mathsf T}A^{-1}B=0$ である。

共有係数を次の有向辺とする。
\[
\begin{array}{c|ccc}
 a&x_4:a\to b&x_3:a\to c&x_1:a\to d\\
 b&x_{13}:b\to a&x_{10}:b\to c&x_6:b\to d\\
 c&x_{18}:c\to a&x_{16}:c\to b&x_7:c\to d\\
 d&x_{19}:d\to a&x_{17}:d\to b&x_{12}:d\to c
\end{array}
\]
対角生成元変換で辺 $i\to j$ の重みは $d_i/d_j$。
全24置換を元の $U,K,A,B,C,J$ へ作用させると、辺は同じ頂点置換で移り、
係数比は $\lambda,\lambda+1$ の単元だけである。
印付けを正規化した後の $\lambda$ の変換と $g$ のアフィン変換も全行列で確認した。
従って非零支持を失わない。
出次数1又は2の頂点を $a$ へ移せば記載の二chartになり、
独立な対角比を選んで非零辺を1にできる。
他の場合は出次数が全て0又は3であり、全零・部分活性・全非零に分かれる。

有限個の支持数は、十二辺の全bitmaskを列挙し、各々へ全24置換を作用させて数える。
各軌道が五場合の一つに含まれること、全4096支持の被覆、
出次数1・2の規格化に必要な置換の存在を元係数の作用と併せて検証した。
これは支持の被覆であり、各支持の全代数係数の分類を意味しない。
\end{proof}

\begin{theorem}\label{thm:cubic-CY-four-point-shared-zero-empty}
四単純基点の零交換子CYにおいて、十二共有係数が全て零となるH3包絡は存在しない。
\end{theorem}
\begin{proof}
$A_*=x_8,B_*=x_{11},C_*=x_{15}$ とする。
元Schur行列の対角六条件は単元倍を除いて
\[
 x_0C_*=x_2B_*=x_5A_*=x_9(B_*-C_*)
 =x_{14}(A_*-C_*)=x_{20}(A_*-B_*)=0.
\]
残余三条件は
\[
\begin{aligned}
 g&=2x_0x_{20}+2C_*(A_*-B_*),\\
 (\lambda+1)/8-g&=2\{x_2x_{14}+B_*(A_*-C_*)\}/\lambda,\\
 -\lambda/8+g&=-2\{x_5x_9+A_*(B_*-C_*)\}/(\lambda+1).
\end{aligned}
\]
これらの線型結合から
$x_0x_{20}-x_2x_{14}+x_5x_9=0$ が従う。

$A_*=B_*=C_*=0$ では、三つの相補積が全非零なら対角尺度変更で
既存 $\mathcal S_{\lambda,t}$ に移る。
一つの相補積が零なら残る二つは非零であり、
全置換と対角規格化で $\mathcal H^0_\lambda$ 又は $\mathcal H^1_\lambda$ に移る。
必要な三つの対角重みの整数rankは3であり、代数閉体上のtorusの像が全体を覆うので、
根を選ぶことによる取りこぼしはない。
命題\ref{prop:cubic-fs-sklyanin}と命題\ref{prop:cubic-fs-remaining}で全排除される。

$(A_*,B_*,C_*)\ne0$ では三つの相補積は全て零である。
実際、一つが非零なら対角条件は例えば $C_*=0,A_*=B_*\ne0$ を強制し、
他の二積は零となるため上の線型恒等式に矛盾する。
全対角係数も零なら量子型 $\mathcal Q$ となり、
\[
 F=A_*B_*+\lambda A_*C_*-(\lambda+1)B_*C_*-\lambda(\lambda+1)/16=0.
\]
対角係数が一つでも非零なら、置換後 $x_0\ne0$ とできる。
すると $C_*=0,A_*B_*=\lambda(\lambda+1)/16\ne0$、
他の五対角係数は全て零であり、$x_0=1$ とすれば $\mathcal O$ に一致する。
$A_*=\pm\lambda/4,\pm(\lambda+1)/4$ の境界には、
係数1の実際の単語積関係が現れてdomain性に反する。
残る全値は命題\ref{prop:cubic-fs-remaining}で排除される。

最後に量子型では、各対の関係は二係数の和が1となる二項式である。
一係数零なら実積関係でdomainに反し、両係数非零ならskew多項式環となる。
$F=0$ は各行の量子係数積が1となる恒等式を与える。
また二係数の和が1なので、全ての $q_{ij}$ は1と異なる。
従って補題\ref{lem:cubic-cy-skew-zero-commutator}により、
中心三次は零又は一本の純三乗であり、これもH3包絡にならない。
全skew領域一般に存在するH3を、この四単純基点の量子型に混同していない。
以上で全枝が尽くされた。
\end{proof}

独立検証器は\\
\texttt{verify\_cubic\_CY\_four\_point\_support.py}、\\
\texttt{verify\_cubic\_CY\_four\_point\_shared\_zero.py}。
元frameの全核、全印付き作用、全支持、共有零の全残余式、
四つの $\mathcal O$ 境界の実積関係、量子型の六二項式と四つのCY行積を検証する。
共有係数非零の全必要係数被覆、全CY切断と相対分類は未完成である。
